\section{Metrics: what we score and why}

\subsection{The four we use}

For actual $y_t$ and forecast $\hat{y}_t$ over $n$ hours:

\begin{align}
  \mathrm{MAE}  &= \tfrac{1}{n}\sum_t |y_t - \hat{y}_t|
  &\mathrm{RMSE} &= \sqrt{\tfrac{1}{n}\sum_t (y_t - \hat{y}_t)^2} \\
  \mathrm{MAPE} &= \tfrac{100}{n}\sum_t \left|\tfrac{y_t - \hat{y}_t}{y_t}\right|
  &\mathrm{PB}_q &= \tfrac{1}{n}\sum_t \max\!\big(q\,d_t,\,(q-1)\,d_t\big),\; d_t = y_t - \hat{y}_t
\end{align}

\subsection{Pros and cons, one by one}

\textbf{MAE (MW).} Average miss in physical units.
Pro: robust, every hour counts equally, easy to explain to a manager.
Con: blind to whether errors cluster into rare disasters.

\textbf{RMSE (MW).} Squares errors before averaging.
Pro: punishes big misses --- and big misses are what cost money on the
balancing market. Con: one broken hour can dominate the score; sensitive to
data glitches.

\textbf{MAPE (\%).} Scale-free.
Pro: comparable across zones and seasons; the industry's headline number.
Con: explodes when $y_t \to 0$ (not an issue for country load, fatal for
prices, which cross zero --- we will NOT use MAPE for prices);
asymmetric --- over-forecasting scores slightly better than under-forecasting.

\textbf{Pinball loss (MW).} The only one that scores a \emph{quantile}.
For P90, under-prediction costs $0.9$ per MW, over-prediction $0.1$.
Pro: rewards honest uncertainty bands; a P90 that is too safe or too tight
both lose. Con: harder to explain; one number per quantile, not one overall.

\subsection{The headline: skill over seasonal naive}

\begin{equation}
  \mathrm{skill} = 1 - \frac{\mathrm{MAE}_{\mathrm{model}}}{\mathrm{MAE}_{\mathrm{naive}}}
\end{equation}

Skill $= 0.25$ means: 25\% less error than copying last week.
Raw MAE flatters everyone in calm weeks; skill cannot be gamed that way,
because the naive baseline enjoys the same calm.

\subsection{Trap to remember}

Never rank models on pooled $R^2$: near-flat periods shrink its denominator
and one pooled number can hide many negative-$R^2$ windows.
RMSE at fixed horizon plus skill-over-naive; per-window breakdowns.

\subsection{Pinball loss: where it comes from}

The name: the error line bends at the quantile like a pinball off a flipper.
Worked example, P90 forecast $\hat{y} = 21{,}000$~MW, quantile $q = 0.9$:

\begin{itemize}
  \item Actual 22,000 (under-predicted, $d = +1000$): loss $= 0.9 \cdot 1000 = 900$.
  \item Actual 20,000 (over-predicted, $d = -1000$): loss $= 0.1 \cdot 1000 = 100$.
\end{itemize}

A well-calibrated P90 sits above the actual $\sim$90\% of the time; the
$9{:}1$ penalty ratio enforces exactly that. Minimizing pinball loss recovers
the true quantile --- it is a \emph{proper scoring rule} for quantiles.

Sources:
\begin{itemize}
  \item Koenker \& Bassett (1978), ``Regression Quantiles'', \emph{Econometrica}
        46(1) --- the original quantile-regression loss.
  \item Gneiting \& Raftery (2007), ``Strictly Proper Scoring Rules, Prediction,
        and Estimation'', \emph{JASA} 102(477) --- why minimizing it gives
        honest quantiles.
  \item Hong et al. (2016), ``Probabilistic energy forecasting: GEFCom2014 and
        beyond'', \emph{Int.\ J.\ Forecasting} 32(3) --- pinball as the official
        score of the largest energy-forecasting competition.
\end{itemize}

\subsection{Interview line}

``MAPE for the manager, RMSE for the balancing-cost risk, pinball for the
quantiles, and skill-over-naive as the headline --- because a model that
cannot beat last week's copy has learned nothing.''
